\chapter{Limites e Continuidade de funções de várias variáveis}
Neste capítulo, introduziremos técnicas referentes a análise da continuidade e cálculo de limites de funções de várias variáveis a valores de $\mathbb R^n$.

Ou seja, vamos estudar funções $f: A\to \mathbb R^m$, onde $A\subseteq \mathbb R^n$ é um conjunto e $m, n\geq 1$ são inteiros.


Ao longo de todo o capítulo, consideraremos que os espaços $\mathbb R^n$ e seus subespaços são munidos da distância euclidiana, genericamente denotada por $d$.

Ou seja, se $x=(x_1, \dots, x_n)$ e $y=(y_1, \dots, y_n)$ são pontos de $\mathbb R^n$, então
\[
    d(x, y) = \sqrt{(x_1 - y_1)^2 + \cdots + (x_n - y_n)^2}.
\]

Além disso, sempre que não especificados, $m, n$ são inteiros positivos.

\section{Limites e continuidade em coordenadas}

Ao estudar funções que saem de um domínio e chegam em $\mathbb R^m$, é muito frequente que possamos reduzir o estudo de suas propriedades ao estudo de suas funções coordenadas.

Lembremos a definição de função coordenada.

\begin{definition}
    Sejam $A$ e $B$ conjuntos e seja $f: A \to B^m$ uma função.

    Para cada $a \in A$, $f(a)$, sendo um elemento de $B^m$, é uma $m$-upla $(b_1, \dots, b_m)$.
    As coordenadas $b_i$ são denotadas por $f_i(a)$.

    Dessa forma, $f_1(a), \dots, f_m(a)$ são os únicos elementos de $B$ tais que
    \[
        f(a) = (f_1(a), \dots, f_m(a)).
    \]

    Para cada $i\in\{1,\dots,m\}$, a função $f_i: A \to B$ assim definida é chamada de \emph{$i$-ésima função coordenada} de $f$.
\end{definition}

\begin{example}
    Seja $f: \mathbb R^2\to \mathbb R^3$ dada por $f(x, y) = (x+y, x-y, xy)$.
    Então as funções coordenadas de $f$ são dadas por:

    \begin{align*}
         & f_1: \mathbb R^2\to \mathbb R, \quad f_1(x, y) = x+y;\\
         & f_2: \mathbb R^2\to \mathbb R, \quad f_2(x, y) = x-y;\\
         & f_3: \mathbb R^2\to \mathbb R, \quad f_3(x, y) = xy.
    \end{align*}
\end{example}

Um caso importante é o das projeções, que são as funções coordenadas da função identidade de $\mathbb R^n$.

\begin{definition}
    Para cada $i\in\{1,\dots,n\}$, a $i$-ésima projeção de $\mathbb R^n$ em $\mathbb R$, denotada por $\pi_i$, é a $i$-ésima função coordenada da função identidade de $\mathbb R^n$.

    Ou seja, $\pi_i: \mathbb R^n\to \mathbb R$ é a função dada por

    \[
        \pi_i(x_1, \dots, x_n) = x_i.
    \]
\end{definition}

\begin{lemma}
    Seja $n\geq 1$ e $i$ um inteiro com $1\leq i\leq n$.

    Então a $i$-ésima projeção de $\mathbb R^n$ em $\mathbb R$ é contínua.

    Ou seja, a função $\pi_i: \mathbb R^n\to \mathbb R$ dada por
    \[
        \pi_i(x_1, \dots, x_n) = x_i
    \]

    é contínua em todos os pontos de $\mathbb R^n$.
\end{lemma}

\begin{proof}
    Seja $p=(p_1, \dots, p_n) \in \mathbb R^n$ arbitrário e $\epsilon>0$.

    Devemos ver que existe $\delta>0$ tal que, para todo $x=(x_1, \dots, x_n) \in \mathbb R^n$, se $d(x, p) < \delta$, então $d(\pi_i(x), \pi_i(p)) < \epsilon$.

    Antes de sugerir um $\delta$, vamos tomar um $\delta>0$ arbitrário e explorar as implicações de $d(x, p) < \delta$.

    Se $x\in \mathbb R^n$ é tal que $d(x, p) < \delta$, então

    \[
        \sqrt{(x_1 - p_1)^2 + \cdots + (x_n - p_n)^2} < \delta.
    \]

    Ora, então,

    \[
        d(\pi_i(x), \pi_i(p))=|x_i - p_i| = \sqrt{(x_i - p_i)^2} \leq \sqrt{(x_1 - p_1)^2 + \cdots + (x_n - p_n)^2} < \delta.
    \]

    Assim, tomando $\delta = \epsilon$, temos que $d(\pi_i(x), \pi_i(p)) < \epsilon$ sempre que $d(x, p) < \delta$.
    Isso atesta a continuidade de $\pi_i$ em $p$.

    Como $p\in \mathbb R^n$ foi arbitrário, concluímos que $\pi_i$ é contínua em todos os pontos de $\mathbb R^n$.
\end{proof}

Note que, se $f:A\to \mathbb R^m$ é uma função, sendo $f_1, \dots, f_m$ suas funções coordenadas, temos que, para cada $i\in\{1,\dots,m\}$,
\[
    f_i = \pi_i \circ f,
\]
em que $\pi_i:\mathbb R^m\to\mathbb R$ é a $i$-ésima projeção de $\mathbb R^m$ em $\mathbb R$.

De fato, se $a\in A$, então $f_i(a)$ é a $i$-ésima coordenada de $f(a)$, ou seja,
\[
    f_i(a)=\pi_i(f(a)).
\]

\begin{proposition}\index{limites!via coordenadas}
    Sejam:
    \begin{itemize}
        \item $(M, d_M)$ um espaço métrico,
        \item $A\subseteq M$ um subespaço,
        \item $f: A \to \mathbb R^m$ uma função,
        \item $p \in M$ um ponto de acumulação de $A$, e
        \item $L=(L_1, \dots, L_m)\in \mathbb R^m$.
    \end{itemize}

    Então $L$ é limite de $f$ em $p$ se, e somente se, para todo $i \in \{1, \dots, m\}$, $L_i$ é limite de $f_i$ em $p$.
\end{proposition}

\begin{proof}
Primeiro, suponha que $L$ é limite de $f$ em $p$.
Dado $i \in \{1, \dots, m\}$, temos que $f_i=\pi_i\circ f$ e que a função $\pi_i: \mathbb R^m\to \mathbb R$ é contínua. Portanto,
\[
    \lim_{x\to p} f_i(x)
    =
    \lim_{x\to p} \pi_i(f(x))
    =
    \pi_i\left(\lim_{x\to p} f(x)\right)
    =
    \pi_i(L)
    =
    L_i.
\]

    Reciprocamente, suponha que para todo $i$, $L_i$ é limite de $f_i$ em $p$.

    Devemos ver que, dado $\epsilon>0$, existe $\delta>0$ tal que, para todo $x \in A\setminus\{p\}$, se $d_M(x, p) < \delta$, então $d(f(x), L) < \epsilon$.
    Fixe $\epsilon>0$ arbitrário.

    Antes de prosseguirmos com $\epsilon$, fixemos um $\bar \epsilon>0$ arbitrário.
    Nós iremos escolher um $\bar \epsilon$ conveniente posteriormente.

    Por hipótese, para cada $i\in \{1, \dots, m\}$, existe $\delta_i>0$ tal que, para todo $x\in A\setminus\{p\}$, se $d_M(x, p) < \delta_i$, então $|f_i(x) - L_i| < \bar \epsilon$.

    Seja $\delta = \min\{\delta_1, \ldots, \delta_m\}$. Então, se $x \in A\setminus\{p\}$ é tal que $d_M(x, p) < \delta$, temos que:

    \begin{align*}
        d(f(x), L) &= \sqrt{(f_1(x) - L_1)^2 + \cdots + (f_m(x) - L_m)^2} \\
        &< \sqrt{\bar \epsilon^2+\dots+\bar \epsilon^2}= \sqrt{m}\bar \epsilon.
    \end{align*}

    Ou seja:
    \begin{equation}\label{eq:limiteViaCoordenadas}
        d(f(x), L) < \sqrt{m}\bar \epsilon.
    \end{equation}

    Assim, apliquemos a construção acima para $\bar \epsilon = \dfrac{\epsilon}{\sqrt{m}}$, obtendo $\delta$ como antes.
    Então, pela desigualdade \eqref{eq:limiteViaCoordenadas}, se $x \in A\setminus\{p\}$ é tal que $d_M(x, p) < \delta$, temos que
    \[
        d(f(x), L) < \sqrt{m}\bar \epsilon = \sqrt{m}\frac{\epsilon}{\sqrt{m}} = \epsilon.
    \]
    Como $\epsilon>0$ foi arbitrário, concluímos que $L$ é limite de $f$ em $p$.
\end{proof}

A Figura~\ref{fig:continuidadeCoord} ilustra a ideia da demonstração acima no caso $m=2$.

\subimport{continuidadeCoord}{imageLimite}


\begin{corollary}\index{continuidade!via coordenadas}
    Sejam:
    \begin{itemize}
        \item $(M, d_M)$ um espaço métrico,
        \item $f: M \to \mathbb R^m$ uma função, e
        \item $p \in M$.
    \end{itemize}

    Sejam $f_1, \dots, f_m$ as funções coordenadas de $f$.

    Então $f$ é contínua em $p$ se, e somente se, para todo $i \in \{1, \dots, m\}$, $f_i$ é contínua em $p$.

    Consequentemente, $f$ é contínua se, e somente se, para todo $i \in \{1, \dots, m\}$, $f_i$ é contínua.
\end{corollary}

\begin{proof}
    Se $p$ é um ponto isolado de $M$, trivialmente $f$ é contínua em $p$ e cada $f_i$ é contínua em $p$, pois a continuidade é automática em pontos isolados.

    Resta considerar o caso em que $p$ é um ponto de acumulação de $M$.
    Nesse caso, pela proposição anterior:

    \begin{align*}
        & f \text{ é contínua em } p \\
        \Leftrightarrow\quad & \lim_{x\to p} f(x) = f(p) \\
        \Leftrightarrow\quad & \forall i \in \{1, \dots, m\}\, \lim_{x\to p} f_i(x) = f_i(p) \\
        \Leftrightarrow\quad & \forall i \in \{1, \dots, m\}\, f_i \text{ é contínua em } p.
    \end{align*}

    A afirmação sobre continuidade de $f$ segue aplicando a equivalência acima a cada ponto $p\in M$.
\end{proof}

\begin{example}
    Seja $f:\mathbb R\to \mathbb R^2$ dada por $f(x)=(x^2, \sin x)$.

    As funções coordenadas de $f$ são $f_1(x)=x^2$ e $f_2(x)=\sin x$.

    Do nosso conhecimento sobre funções de uma variável, sabemos que $f_1$ e $f_2$ são funções contínuas.
    Assim, $f$ é contínua.
\end{example}

\begin{example}
    Seja $g:\mathbb R\to \mathbb R^2$ dada por

    \[
        g(x) = \begin{cases}
            (x^2, \sin x), & \text{se } x\neq 0;\\
            (0, 1), & \text{se } x=0.
        \end{cases}
    \]

    As funções coordenadas de $g$ são $g_1:\mathbb R\to \mathbb R$ e $g_2:\mathbb R\to \mathbb R$ dadas por

    \[
        g_1(x) = \begin{cases}
            x^2, & \text{se } x\neq 0;\\
            0, & \text{se } x=0,
        \end{cases}
    \]

    \[
        g_2(x) = \begin{cases}
            \sin x, & \text{se } x\neq 0;\\
            1, & \text{se } x=0.
        \end{cases}
    \]

    Note que $g_1(x)=x^2$ para todo $x\in \mathbb R$, e, portanto, $g_1$ é contínua.

    Por outro lado, o conjunto de pontos de continuidade de $g_2$ é $\mathbb R\setminus\{0\}$.

    Portanto, o conjunto de pontos de continuidade de $g$ é a interseção do conjunto de pontos de continuidade de $g_1$ com o conjunto de pontos de continuidade de $g_2$, ou seja, $\mathbb R\setminus\{0\}$.
\end{example}

\begin{example}
    Seja $f:\mathbb R^2\rightarrow \mathbb R^2$ dada por

    \[
        f(x, y)=(\cos(x), \sin(y)).
    \]
    
    A função $f$ é contínua.
    De fato, as funções coordenadas de $f$ são as funções $f_1,f_2:\mathbb R^2\to \mathbb R$ dadas por
    \[
        f_1(x, y)=\cos(x), \quad f_2(x, y)=\sin(y).
    \]

    Temos que $f_1$ é contínua, pois $f_1=\cos\circ \pi_1$, composição de funções contínuas.
    Analogamente, $f_2$ é contínua, pois $f_2=\sin\circ \pi_2$, composição de funções contínuas.
    Portanto, $f$ é contínua.
\end{example}

\section*{Exercícios}

\begin{exercise}
    Em cada item abaixo, justifique a continuidade da função indicada.

    \begin{enumerate}[label=(\alph*)]
        \item $f:\mathbb R^2\to\mathbb R^3$ dada por
        \[
            f(x,y)=(x,y,0).
        \]

        \item $g:\mathbb R\to\mathbb R^2$ dada por
        \[
            g(x)=(x^2,\sin x).
        \]

        \item $h:[0,+\infty[\to\mathbb R^2$ dada por
        \[
            h(x)=(\sqrt{x},x).
        \]

        \item $\varphi:]0,+\infty[^2\to\mathbb R^3$ dada por
        \[
            \varphi(x,y)=\left(\sqrt{x},\frac{1}{y},\sin x\right).
        \]
    \end{enumerate}
\end{exercise}

\begin{exercise}
    Seja $g:\mathbb R\to\mathbb R^2$ dada por
    \[
        g(x)=
        \begin{cases}
            (x,\sin x), & \text{se } x\neq 0,\\
            (0,1), & \text{se } x=0.
        \end{cases}
    \]

    Determine o conjunto dos pontos de continuidade de $g$.
\end{exercise}

\begin{exercise}
    Seja $A=\mathbb R^2\setminus\{(1,2)\}$ e seja
    $f:A\to\mathbb R^3$ dada por
    \[
        f(x,y)=(y,x,5).
    \]

    Calcule justificando com os resultados já apresentados.
    \[
        \lim_{(x,y)\to(1,2)} f(x,y).
    \]
\end{exercise}
\section{Operações com limites e continuidade}

Na seção anterior, vimos que uma função a valores em $\mathbb R^m$ pode ser estudada coordenada a coordenada.
Agora veremos que as operações algébricas usuais também se comportam bem com limites e continuidade.

A ideia central desta seção é simples: se duas funções se aproximam de certos valores, então a soma se aproxima da soma desses valores.
Quando as funções são reais, o mesmo vale para produtos, inversos e quocientes, desde que as expressões estejam definidas.

\begin{proposition}\index{limite!soma}\index{limite!diferença}\index{limite!multiplicação por escalar}
    Sejam:
    \begin{itemize}
        \item $(M,d_M)$ um espaço métrico,
        \item $A\subseteq M$ um subespaço,
        \item $f,g:A\to\mathbb R^m$ funções,
        \item $p\in M$ um ponto de acumulação de $A$,
        \item $L,S\in\mathbb R^m$, e
        \item $\lambda\in\mathbb R$.
    \end{itemize}

    Se
    \[
        \lim_{x\to p} f(x)=L
        \qquad\text{e}\qquad
        \lim_{x\to p} g(x)=S,
    \]
    então:
    \begin{enumerate}[label=(\alph*)]
        \item $\displaystyle \lim_{x\to p} (f+g)(x)=L+S$;
        \item $\displaystyle \lim_{x\to p} (f-g)(x)=L-S$;
        \item $\displaystyle \lim_{x\to p} (\lambda f)(x)=\lambda L$.
    \end{enumerate}
\end{proposition}

\begin{proof}
    Comecemos pela soma.
    Fixe $\epsilon>0$.

    Como $L$ é limite de $f$ em $p$, existe $\delta_f>0$ tal que, para todo $x\in A\setminus\{p\}$,
    \[
        \text{se }d_M(x,p)<\delta_f,\text{ então }d(f(x),L)<\frac{\epsilon}{2}.
    \]

    Analogamente, como $S$ é limite de $g$ em $p$, existe $\delta_g>0$ tal que, para todo $x\in A\setminus\{p\}$,
    \[
        \text{se }d_M(x,p)<\delta_g,\text{ então }d(g(x),S)<\frac{\epsilon}{2}.
    \]

    Tome $\delta=\min\{\delta_f,\delta_g\}$.
    Se $x\in A\setminus\{p\}$ e $d_M(x,p)<\delta$, então:
    \begin{align*}
        d((f+g)(x),L+S)
        &=\|(f(x)+g(x))-(L+S)\|\\
        &=\|(f(x)-L)+(g(x)-S)\|\\
        &\leq \|f(x)-L\|+\|g(x)-S\|\\
        &=d(f(x),L)+d(g(x),S)\\
        &<\frac{\epsilon}{2}+\frac{\epsilon}{2}
        =\epsilon.
    \end{align*}
    Logo, $\lim_{x\to p}(f+g)(x)=L+S$.

    Para a multiplicação por escalar, se $\lambda=0$, então $\lambda f$ é a função constante igual a $0$, e o limite é $0=\lambda L$.

    Suponha, então, que $\lambda\neq 0$.
    Fixe $\epsilon>0$.
    Como $L$ é limite de $f$ em $p$, existe $\delta>0$ tal que, para todo $x\in A\setminus\{p\}$,
    \[
        \text{se }d_M(x,p)<\delta,\text{ então }d(f(x),L)<\frac{\epsilon}{|\lambda|}.
    \]
    Assim, se $x\in A\setminus\{p\}$ e $d_M(x,p)<\delta$, então
    \[
        d((\lambda f)(x),\lambda L)
        =
        \|\lambda f(x)-\lambda L\|
        =
        |\lambda|\|f(x)-L\|
        <
        |\lambda|\frac{\epsilon}{|\lambda|}
        =
        \epsilon.
    \]
    Portanto, $\lim_{x\to p}(\lambda f)(x)=\lambda L$.

    Por fim, a diferença é consequência dos dois itens já provados, pois
    \[
        f-g=f+(-1)g.
    \]
    Assim,
    \[
        \lim_{x\to p}(f-g)(x)
        =
        L+(-1)S
        =
        L-S.
    \]
\end{proof}

\begin{corollary}\index{continuidade!soma}\index{continuidade!diferença}\index{continuidade!multiplicação por escalar}
    Sejam $A\subseteq M$ um subespaço de um espaço métrico $(M,d_M)$ e sejam $f,g:A\to\mathbb R^m$ funções.

    Se $f$ e $g$ são contínuas em $p\in A$, então $f+g$ e $f-g$ são contínuas em $p$.

    Além disso, para todo $\lambda\in\mathbb R$, se $f$ é contínua em $p$, então $\lambda f$ é contínua em $p$.

    Consequentemente, se $f$ e $g$ são contínuas, então $f+g$ e $f-g$ são contínuas; e, para todo $\lambda\in\mathbb R$, $\lambda f$ é contínua.
\end{corollary}

\begin{proof}
    Se $p$ é ponto isolado de $A$, todas as funções com domínio $A$ são contínuas em $p$.

    Suponha, então, que $p$ é ponto de acumulação de $A$.
    Como $f$ e $g$ são contínuas em $p$, temos
    \[
        \lim_{x\to p} f(x)=f(p)
        \qquad\text{e}\qquad
        \lim_{x\to p} g(x)=g(p).
    \]
    Pela proposição anterior,
    \[
        \lim_{x\to p}(f+g)(x)=f(p)+g(p)=(f+g)(p),
    \]
    \[
        \lim_{x\to p}(f-g)(x)=f(p)-g(p)=(f-g)(p),
    \]
    e
    \[
        \lim_{x\to p}(\lambda f)(x)=\lambda f(p)=(\lambda f)(p).
    \]
    Logo, $f+g$, $f-g$ e $\lambda f$ são contínuas em $p$.

    A afirmação global segue aplicando o resultado ponto a ponto.
\end{proof}

\begin{example}
    Seja $F:\mathbb R^2\to\mathbb R^2$ dada por
    \[
        F(x,y)=(x,y)+(y,x).
    \]
    A primeira parcela é a função identidade de $\mathbb R^2$, e a segunda parcela é a função $G(x,y)=(y,x)$, cujas coordenadas são projeções.
    Portanto, ambas são contínuas.
    Assim, $F$ é contínua.

    Naturalmente, simplificando a expressão, temos
    \[
        F(x,y)=(x+y,x+y).
    \]
\end{example}

\begin{lemma}\index{limite!limitação local}
    Sejam:
    \begin{itemize}
        \item $(M,d_M)$ um espaço métrico,
        \item $A\subseteq M$ um subespaço,
        \item $f:A\to\mathbb R$ uma função,
        \item $p\in M$ um ponto de acumulação de $A$, e
        \item $L\in\mathbb R$.
    \end{itemize}

    Se $\lim_{x\to p}f(x)=L$, então $f$ é limitada em alguma vizinhança perfurada de $p$.

    Mais precisamente, existem $\rho>0$ e $C>0$ tais que, para todo $x\in A\setminus\{p\}$,
    \[
        d_M(x,p)<\rho
        \quad\Longrightarrow\quad
        |f(x)|\leq C.
    \]
\end{lemma}

\begin{proof}
    Apliquemos a definição de limite com margem de erro igual a $1$.
    Existe $\rho>0$ tal que, para todo $x\in A\setminus\{p\}$,
    \[
        \text{se }d_M(x,p)<\rho,\text{ então } |f(x)-L|<1.
    \]

    Para tais $x$, temos:
    \[
        |f(x)|
        =
        |f(x)-L+L|
        \leq
        |f(x)-L|+|L|
        <
        1+|L|.
    \]

    Portanto, podemos tomar $C=1+|L|$.
\end{proof}

\begin{proposition}\index{limite!produto}
    Sejam:
    \begin{itemize}
        \item $(M,d_M)$ um espaço métrico,
        \item $A\subseteq M$ um subespaço,
        \item $f,g:A\to\mathbb R$ funções,
        \item $p\in M$ um ponto de acumulação de $A$, e
        \item $L,S\in\mathbb R$.
    \end{itemize}

    Se
    \[
        \lim_{x\to p}f(x)=L
        \qquad\text{e}\qquad
        \lim_{x\to p}g(x)=S,
    \]
    então
    \[
        \lim_{x\to p}(fg)(x)=LS.
    \]
\end{proposition}

\begin{proof}
    Pelo lema anterior, existem $\rho>0$ e $C>0$ tais que, para todo $x\in A\setminus\{p\}$,
    \[
        d_M(x,p)<\rho
        \quad\Longrightarrow\quad
        |f(x)|\leq C.
    \]

    Fixe $\epsilon>0$.

    Como $L$ é limite de $f$ em $p$, existe $\delta_f>0$ tal que, para todo $x\in A\setminus\{p\}$,
    \[
        \text{se }d_M(x,p)<\delta_f,\text{ então } |f(x)-L|<\frac{\epsilon}{2(|S|+1)}.
    \]

    Como $S$ é limite de $g$ em $p$, existe $\delta_g>0$ tal que, para todo $x\in A\setminus\{p\}$,
    \[
        \text{se }d_M(x,p)<\delta_g,\text{ então } |g(x)-S|<\frac{\epsilon}{2C}.
    \]

    Tome $\delta=\min\{\rho,\delta_f,\delta_g\}$.
    Se $x\in A\setminus\{p\}$ e $d_M(x,p)<\delta$, então:
    \begin{align*}
        |(fg)(x)-LS|
        &=
        |f(x)g(x)-LS|\\
        &=
        |f(x)g(x)-f(x)S+f(x)S-LS|\\
        &=
        |f(x)(g(x)-S)+S(f(x)-L)|\\
        &\leq
        |f(x)|\,|g(x)-S|+|S|\,|f(x)-L|\\
        &<
        C\frac{\epsilon}{2C}
        +
        |S|\frac{\epsilon}{2(|S|+1)}\\
        &<
        \frac{\epsilon}{2}+\frac{\epsilon}{2}
        =
        \epsilon.
    \end{align*}

    Logo, $\lim_{x\to p}(fg)(x)=LS$.
\end{proof}

\begin{corollary}\index{continuidade!produto}
    Sejam $A\subseteq M$ um subespaço de um espaço métrico $(M,d_M)$ e sejam $f,g:A\to\mathbb R$ funções.

    Se $f$ e $g$ são contínuas em $p\in A$, então $fg$ é contínua em $p$.

    Consequentemente, se $f$ e $g$ são contínuas, então $fg$ é contínua.
\end{corollary}

\begin{proof}
    Se $p$ é ponto isolado de $A$, então $fg$ é contínua em $p$.

    Suponha que $p$ é ponto de acumulação de $A$.
    Como $f$ e $g$ são contínuas em $p$, temos
    \[
        \lim_{x\to p}f(x)=f(p)
        \qquad\text{e}\qquad
        \lim_{x\to p}g(x)=g(p).
    \]
    Pela proposição anterior,
    \[
        \lim_{x\to p}(fg)(x)=f(p)g(p)=(fg)(p).
    \]
    Portanto, $fg$ é contínua em $p$.

    A afirmação global segue aplicando o resultado em cada ponto do domínio.
\end{proof}

\begin{example}
    A função $P:\mathbb R^3\to\mathbb R$ dada por
    \[
        P(x,y,z)=3x^2yz-7xz+5
    \]
    é contínua.

    De fato, as funções $(x,y,z)\mapsto x$, $(x,y,z)\mapsto y$ e $(x,y,z)\mapsto z$ são projeções, portanto contínuas.
    Produtos, multiplicações por escalares e somas preservam continuidade.
    Logo, $P$ é contínua.
\end{example}

\begin{proposition}\index{limite!inversa multiplicativa}
    Sejam:
    \begin{itemize}
        \item $(M,d_M)$ um espaço métrico,
        \item $A\subseteq M$ um subespaço,
        \item $f:A\to\mathbb R$ uma função,
        \item $p\in M$ um ponto de acumulação de $A$, e
        \item $L\in\mathbb R\setminus\{0\}$.
    \end{itemize}

    Suponha que
    \[
        \lim_{x\to p}f(x)=L.
    \]

    Seja
    \[
        B=\{x\in A:f(x)\neq 0\}.
    \]

    Então $p$ é ponto de acumulação de $B$ e
    \[
        \lim_{x\to p}\frac{1}{f(x)}=\frac{1}{L},
    \]
    onde $\frac{1}{f}$ é vista como função de $B$ em $\mathbb R$.
\end{proposition}

\begin{proof}
    Como $L\neq 0$, temos $\frac{|L|}{2}>0$.
    Pela hipótese de limite, existe $\rho>0$ tal que, para todo $x\in A\setminus\{p\}$,
    \[
        \text{se }d_M(x,p)<\rho,\text{ então } |f(x)-L|<\frac{|L|}{2}.
    \]

    Para tais $x$, temos
    \[
        |f(x)|
        =
        |L+(f(x)-L)|
        \geq
        |L|-|f(x)-L|
        >
        |L|-\frac{|L|}{2}
        =
        \frac{|L|}{2}.
    \]
    Em particular, $f(x)\neq 0$.
    Portanto, todo ponto de $A\setminus\{p\}$ suficientemente próximo de $p$ pertence a $B$.
    Como $p$ é ponto de acumulação de $A$, segue que $p$ também é ponto de acumulação de $B$.

    Fixe agora $\epsilon>0$.
    Como $L$ é limite de $f$ em $p$, existe $\delta_1>0$ tal que, para todo $x\in A\setminus\{p\}$,
    \[
        \text{se }d_M(x,p)<\delta_1,\text{ então }
        |f(x)-L|<\frac{\epsilon |L|^2}{2}.
    \]

    Tome $\delta=\min\{\rho,\delta_1\}$.
    Se $x\in B\setminus\{p\}$ e $d_M(x,p)<\delta$, então $x\in A\setminus\{p\}$, $|f(x)|>\frac{|L|}{2}$ e:
    \[
        \left|\frac{1}{f(x)}-\frac{1}{L}\right|
        =
        \left|\frac{L-f(x)}{f(x)L}\right|
        =
        \frac{|f(x)-L|}{|f(x)|\,|L|}
        <
        \frac{|f(x)-L|}{\frac{|L|}{2}|L|}
        =
        \frac{2|f(x)-L|}{|L|^2}
        <
        \epsilon.
    \]

    Logo,
    \[
        \lim_{x\to p}\frac{1}{f(x)}=\frac{1}{L}.
    \]
\end{proof}

\begin{corollary}\index{limite!quociente}
    Sejam:
    \begin{itemize}
        \item $(M,d_M)$ um espaço métrico,
        \item $A\subseteq M$ um subespaço,
        \item $f,g:A\to\mathbb R$ funções,
        \item $p\in M$ um ponto de acumulação de $A$, e
        \item $L,S\in\mathbb R$, com $S\neq 0$.
    \end{itemize}

    Se
    \[
        \lim_{x\to p}f(x)=L
        \qquad\text{e}\qquad
        \lim_{x\to p}g(x)=S,
    \]
    então, pondo
    \[
        B=\{x\in A:g(x)\neq 0\},
    \]
    temos que $p$ é ponto de acumulação de $B$ e
    \[
        \lim_{x\to p}\frac{f(x)}{g(x)}=\frac{L}{S},
    \]
    onde $\frac{f}{g}$ é vista como função de $B$ em $\mathbb R$.
\end{corollary}

\begin{proof}
    Pela proposição anterior aplicada a $g$, temos que $p$ é ponto de acumulação de $B$ e
    \[
        \lim_{x\to p}\frac{1}{g(x)}=\frac{1}{S}.
    \]

    Além disso, como $B\subseteq A$, a mesma escolha de $\delta$ que mostra que $f(x)$ se aproxima de $L$ para $x\in A\setminus\{p\}$ também mostra que $f(x)$ se aproxima de $L$ para $x\in B\setminus\{p\}$.
    Assim,
    \[
        \lim_{x\to p} f(x)=L
    \]
    quando $f$ é considerada apenas no domínio $B$.

    Aplicando a regra do produto no domínio $B$, obtemos:
    \[
        \lim_{x\to p}\frac{f(x)}{g(x)}
        =
        \lim_{x\to p}\left(f(x)\frac{1}{g(x)}\right)
        =
        L\frac{1}{S}
        =
        \frac{L}{S}.
    \]
\end{proof}

\begin{corollary}\index{continuidade!inversa multiplicativa}\index{continuidade!quociente}
    Sejam $A\subseteq M$ um subespaço de um espaço métrico $(M,d_M)$ e sejam $f,g:A\to\mathbb R$ funções.

    \begin{enumerate}[label=(\alph*)]
        \item Se $f$ é contínua em $p\in A$ e $f(p)\neq 0$, então $\frac{1}{f}$ é contínua em $p$, quando considerada no domínio $\{x\in A:f(x)\neq 0\}$.
        \item Se $f$ e $g$ são contínuas em $p\in A$ e $g(p)\neq 0$, então $\frac{f}{g}$ é contínua em $p$, quando considerada no domínio $\{x\in A:g(x)\neq 0\}$.
    \end{enumerate}

    Consequentemente, se $g(x)\neq 0$ para todo $x\in A$ e $f,g$ são contínuas, então $\frac{f}{g}:A\to\mathbb R$ é contínua.
\end{corollary}

\begin{proof}
    Provemos o item (a).
    Seja $B=\{x\in A:f(x)\neq 0\}$.
    Como $f(p)\neq 0$, temos $p\in B$.

    Se $p$ é ponto isolado de $B$, então $\frac{1}{f}$ é contínua em $p$.

    Suponha que $p$ é ponto de acumulação de $B$.
    Como $f$ é contínua em $p$, temos
    \[
        \lim_{x\to p} f(x)=f(p).
    \]
    Pelo resultado sobre inversas, considerando $f$ restrita a $B$, temos
    \[
        \lim_{x\to p}\frac{1}{f(x)}=\frac{1}{f(p)}.
    \]
    Como $\frac{1}{f}(p)=\frac{1}{f(p)}$, segue que $\frac{1}{f}$ é contínua em $p$.

    O item (b) segue do item (a) e da continuidade de produtos:
    \[
        \frac{f}{g}=f\frac{1}{g}.
    \]

    A afirmação global segue aplicando o item (b) em cada ponto de $A$.
\end{proof}

\begin{example}
    Seja $R:\mathbb R^2\setminus\{(x,y):x^2+y^2=1\}\to\mathbb R$ dada por
    \[
        R(x,y)=\frac{x^2+y}{x^2+y^2-1}.
    \]

    O numerador e o denominador são funções polinomiais, logo contínuas.
    No domínio indicado, o denominador nunca se anula.
    Portanto, $R$ é contínua.
\end{example}

\begin{corollary}\index{função polinomial}\index{continuidade!função polinomial}
    Toda função polinomial $P:\mathbb R^n\to\mathbb R$ é contínua.
\end{corollary}

\begin{proof}
    Uma função polinomial é obtida a partir de funções constantes e projeções por um número finito de somas, produtos e multiplicações por escalares.
    Funções constantes e projeções são contínuas.
    Como as operações acima preservam continuidade, toda função polinomial é contínua.
\end{proof}

\begin{corollary}\index{função racional}\index{continuidade!função racional}
    Se $P,Q:\mathbb R^n\to\mathbb R$ são funções polinomiais e
    \[
        D=\{x\in\mathbb R^n:Q(x)\neq 0\},
    \]
    então a função racional $\frac{P}{Q}:D\to\mathbb R$ é contínua.
\end{corollary}

\begin{proof}
    Pelo corolário anterior, $P$ e $Q$ são contínuas.
    No domínio $D$, o denominador $Q$ não se anula.
    Portanto, o quociente $\frac{P}{Q}$ é contínuo em $D$.
\end{proof}

\section*{Exercícios}

\begin{exercise}
    Calcule os limites abaixo, justificando com as regras de operações.
    \begin{enumerate}[label=(\alph*)]
        \item $\displaystyle \lim_{(x,y)\to(1,2)}(x^2y+3x-4y)$.
        \item $\displaystyle \lim_{(x,y,z)\to(0,1,-1)}(xz+y^2)$.
        \item $\displaystyle \lim_{(x,y)\to(2,1)}\frac{x^2+y}{x-y}$.
    \end{enumerate}
\end{exercise}

\begin{exercise}
    Seja $f:\mathbb R^2\to\mathbb R^3$ dada por
    \[
        f(x,y)=(x^2+y,\ xy,\ 3x-y).
    \]
    Mostre que $f$ é contínua.
\end{exercise}

\begin{exercise}
    Determine o conjunto dos pontos de continuidade da função $f:\mathbb R^2\setminus\{(0,0)\}\to\mathbb R$ dada por
    \[
        f(x,y)=\frac{x-y}{x^2+y^2}.
    \]
\end{exercise}

\begin{exercise}
    Seja $A\subseteq M$ um subespaço de um espaço métrico e sejam $f,g:A\to\mathbb R$ funções contínuas.
    Mostre que a função $h:A\to\mathbb R^2$ dada por
    \[
        h(x)=(f(x)+g(x), f(x)g(x))
    \]
    é contínua.
\end{exercise}

\begin{exercise}
    Seja $f:A\to\mathbb R$ uma função e seja $p$ um ponto de acumulação de $A$.
    Suponha que $\lim_{x\to p}f(x)=L$.
    Mostre diretamente da definição que
    \[
        \lim_{x\to p}|f(x)|=|L|.
    \]
\end{exercise}

\section{O Teorema do Confronto}\index{Teorema do Confronto}\index{sanduíche, teorema do|see{Teorema do Confronto}}


\begin{proposition}Sejam
    \begin{itemize}
        \item $A\subseteq \mathbb R^n$ um conjunto,
        \item $f, g: A \to \mathbb R$ funções,
        \item $p \in \mathbb R^n$ um ponto de acumulação de $A$,
        \item $L=\lim_{x\to p} f(x)$, e,
        \item $M=\lim_{x\to p} g(x)$.
    \end{itemize}

    Então:
    \begin{enumerate}[label=(\alph*)]
        \item $\displaystyle \lim_{x\to p}(f+g)(x) = L + M$,
        \item $\displaystyle \lim_{x\to p}(f-g)(x) = L - M$, e 
        \item $\displaystyle \lim_{x\to p}(fg)(x) = LM$.
    \end{enumerate}
\end{proposition}

\begin{proposition}Sejam
    \begin{itemize}
        \item $A\subseteq \mathbb R^n$ um conjunto,
        \item $f, g: A \to \mathbb R$ funções,
        \item $p \in A$.
    \end{itemize}

    Se $f$ e $g$ são contínuas em $p$, então:
    \begin{enumerate}[label=(\alph*)]
        \item $f+g$ é contínua em $p$,
        \item $f-g$ é contínua em $p$, e
        \item $fg$ é contínua em $p$.
    \end{enumerate}

    Consequentemente, se $f$ e $g$ são contínuas, então $f+g$, $f-g$ e $fg$ são contínuas.
\end{proposition}

\begin{proposition}Sejam
    \begin{itemize}
        \item $A\subseteq \mathbb R^n$ um conjunto,
        \item $f, g: A \to \mathbb R$ funções com $g(x)\neq 0$ para todo $x \in A$,
        \item $p \in \mathbb R^n$ um ponto de acumulação de $A$,
        \item $L=\lim_{x\to p} f(x)$, e,
        \item $M=\lim_{x\to p} g(x)$.
    \end{itemize}

    Então:
    \[
        \lim_{x\to p}\frac{f(x)}{g(x)} = \frac{L}{M}.
    \]
\end{proposition}

\begin{proposition}Sejam
    \begin{itemize}
        \item $A\subseteq \mathbb R^n$ um conjunto,
        \item $f, g: A \to \mathbb R$ funções com $g(x)\neq 0$ para todo $x \in A$,
        \item $p\in A$.
    \end{itemize}

    Se $f$ e $g$ são contínuas em $p$, então $\displaystyle\frac{f}{g}$ é contínua em $p$.

    Consequentemente, se $f$ e $g$ são contínuas, então $\displaystyle\frac{f}{g}$ é contínua.
\end{proposition}


\begin{proposition}\index{limite!Teorema do Confronto}
    Sejam $f, g, h: A \to \mathbb R$ uma função, onde $A\subseteq \mathbb R^m$.

    Sejam $p \in A$ um ponto de acumulação de $A$ e $L \in \mathbb R$.

    Suponha que existe uma bola aberta $B$ em torno de $p$ tal que, para todo $x \in B\cap A\setminus\{p\}$, temos que $f(x) \leq g(x) \leq h(x)$.

    Nessas hipóteses, se $\lim_{x\to p} f(x) = \lim_{x\to p} h(x) = L$, então $\lim_{x\to p} g(x) = L$.
\end{proposition}

\begin{proposition}\index{limite!Teorema do Confronto}
    Sejam
    \begin{itemize}
        \item $A\subseteq \mathbb R^n$ um conjunto,
        \item $f, g, h: A \to \mathbb R$ funções,
        \item $p \in \mathbb R^n$ um ponto de acumulação de $A$, e
        \item $L \in \mathbb R$.
    \end{itemize}

    Suponha que existe uma bola aberta $B$ em torno de $p$ tal que, para todo $x \in B\cap A\setminus\{p\}$, temos
    \[
        f(x) \leq g(x) \leq h(x).
    \]

    Então:
    \[
        \text{se } \lim_{x\to p} f(x) = \lim_{x\to p} h(x) = L, \text{ então } \lim_{x\to p} g(x) = L.
    \]
\end{proposition}

\begin{proof}
    Seja dado $\epsilon > 0$.
    
    Como $\lim_{x\to p} f(x) = L$, existe $\delta_f>0$ tal que, para todo $x \in A\setminus\{p\}$, se $d(x, p) < \delta_f$, então $|f(x) - L| < \epsilon$.

    Analogamente, como $\lim_{x\to p} h(x) = L$, existe $\delta_h>0$ tal que, para todo $x \in A\setminus\{p\}$, se $d(x, p) < \delta_h$, então $|h(x) - L| < \epsilon$.

    Por fim, existe $\delta_g>0$ tal que, para todo $x \in B(p, \delta_g)$, se $x \in A\setminus\{p\}$, então $f(x) \leq g(x) \leq h(x)$.

    Seja $\delta = \min\{\delta_f, \delta_h, \delta_g\}$.

    Seja $x \in A\setminus\{p\}$ tal que $d(x, p) < \delta$, então $-\epsilon<f(x) - L \leq g(x) - L \leq h(x) - L < \epsilon$, logo, $|g(x) - L| < \epsilon$, ou seja, $d(g(x), L) < \epsilon$.
\end{proof}
\begin{corollary}
    Sejam
    \begin{itemize}
        \item $A\subseteq \mathbb R^n$ um conjunto,
        \item $f, g$ funções,
        \item $p \in \mathbb R^n$ um ponto de acumulação de $A$, e
        \item $L \in \mathbb R$.
    \end{itemize}

    Suponha que:
    \begin{itemize}
        \item existe uma bola aberta $B$ em torno de $p$ tal que $g|_B$ é limitada, ou seja, tal que existe $M>0$ tal que, para todo $x \in B\cap A\setminus\{p\}$, temos $|g(x)| < M$, e
        \item $\lim_{x\to p} f(x) = 0$.
     \end{itemize}

     Então $\lim_{x\to p} f(x)g(x) = 0$.
\end{corollary}

\begin{proposition}Sejam
    \begin{itemize}
        \item $I\subseteq \mathbb R$ um intervalo aberto
        \item $A\subseteq \mathbb R^n$ um conjunto,
        \item $a \in \mathbb R^n$ um ponto de acumulação de $A$,
        \item $\gamma: I \to \mathbb R^n$ uma função
        \item $t_0 \in I$ com $\gamma(t_0) = a$
    \end{itemize}

    Assuma que:
    \begin{itemize}
        \item $\gamma(t)\neq a$ para todo $t \in I\setminus\{t_0\}$,
        \item $\lim_{t\to t_0} \gamma(t) = a$, e
        \item $\lim_{t\to t_0} f(\gamma(t)) = L$.
    \end{itemize}

    Então $\lim_{t\to t_0} f(\gamma(t)) = L$.
\end{proposition}
\section{Testes de continuidade via curvas}\index{teste das curvas}\index{curvas!teste de continuidade}

Funções de várias variáveis podem ser difíceis de visualizar diretamente.
Uma técnica muito útil é aproximar o ponto por curvas.
Ao compor uma função de várias variáveis com uma curva, obtemos uma função de uma variável, que costuma ser mais simples de estudar.

O princípio básico é o seguinte: se uma função possui limite em um ponto, então esse mesmo limite aparece ao longo de qualquer curva que se aproxima desse ponto sem ficar presa nele.

\begin{proposition}\index{limite!ao longo de curvas}
    Sejam:
    \begin{itemize}
        \item $A\subseteq\mathbb R^n$ um conjunto,
        \item $f:A\to\mathbb R^m$ uma função,
        \item $p\in\mathbb R^n$ um ponto de acumulação de $A$,
        \item $L\in\mathbb R^m$,
        \item $I\subseteq\mathbb R$ um conjunto,
        \item $t_0\in\mathbb R$ um ponto de acumulação de $I$, e
        \item $\gamma:I\to A$ uma função.
    \end{itemize}

    Suponha que:
    \begin{itemize}
        \item $\lim_{t\to t_0}\gamma(t)=p$;
        \item existe $\eta>0$ tal que, para todo $t\in I\setminus\{t_0\}$,
        \[
            |t-t_0|<\eta
            \quad\Longrightarrow\quad
            \gamma(t)\neq p;
        \]
        \item $\lim_{x\to p}f(x)=L$.
    \end{itemize}

    Então
    \[
        \lim_{t\to t_0}f(\gamma(t))=L.
    \]
\end{proposition}

\begin{proof}
    Fixe $\epsilon>0$.

    Como $\lim_{x\to p}f(x)=L$, existe $\delta_f>0$ tal que, para todo $x\in A\setminus\{p\}$,
    \[
        \text{se }d(x,p)<\delta_f,\text{ então }d(f(x),L)<\epsilon.
    \]

    Como $\lim_{t\to t_0}\gamma(t)=p$, existe $\delta_\gamma>0$ tal que, para todo $t\in I\setminus\{t_0\}$,
    \[
        \text{se }|t-t_0|<\delta_\gamma,\text{ então }d(\gamma(t),p)<\delta_f.
    \]

    Tome
    \[
        \delta=\min\{\eta,\delta_\gamma\}.
    \]

    Se $t\in I\setminus\{t_0\}$ e $|t-t_0|<\delta$, então $\gamma(t)\neq p$ e $d(\gamma(t),p)<\delta_f$.
    Como $\gamma(t)\in A$, temos $\gamma(t)\in A\setminus\{p\}$.
    Portanto,
    \[
        d(f(\gamma(t)),L)<\epsilon.
    \]

    Logo,
    \[
        \lim_{t\to t_0}f(\gamma(t))=L.
    \]
\end{proof}

\begin{corollary}[Teste negativo por curvas]\index{limite!teste negativo por curvas}
    Seja $f:A\to\mathbb R^m$, com $A\subseteq\mathbb R^n$, e seja $p$ um ponto de acumulação de $A$.

    Se existir uma curva $\gamma:I\to A$ tal que $\gamma(t)\to p$ quando $t\to t_0$, $\gamma(t)\neq p$ para $t\neq t_0$ suficientemente próximo de $t_0$, e $f\circ\gamma$ não possuir limite em $t_0$, então $f$ não possui limite em $p$.

    Além disso, se existirem duas curvas $\gamma_1:I_1\to A$ e $\gamma_2:I_2\to A$ que se aproximam de $p$ e tais que
    \[
        \lim_{t\to t_1}f(\gamma_1(t))
        \neq
        \lim_{t\to t_2}f(\gamma_2(t)),
    \]
    então $f$ não possui limite em $p$.
\end{corollary}

\begin{proof}
    Se $f$ possuísse limite em $p$, então, pela proposição anterior, toda curva que se aproxima de $p$ satisfazendo as hipóteses teria, após composição com $f$, esse mesmo limite.
    Portanto, uma curva sem limite, ou duas curvas com limites diferentes, impossibilitam a existência do limite de $f$ em $p$.
\end{proof}

\begin{corollary}\index{continuidade!teste negativo por curvas}
    Seja $A\subseteq\mathbb R^n$, seja $f:A\to\mathbb R^m$ uma função e seja $p\in A$.

    Se $f$ é contínua em $p$, então, para toda curva contínua $\gamma:I\to A$ e todo $t_0\in I$ tal que $\gamma(t_0)=p$, a função
    \[
        f\circ\gamma:I\to\mathbb R^m
    \]
    é contínua em $t_0$.

    Consequentemente, se existe uma curva contínua $\gamma:I\to A$ com $\gamma(t_0)=p$ tal que $f\circ\gamma$ não é contínua em $t_0$, então $f$ não é contínua em $p$.
\end{corollary}

\begin{proof}
    Fixe $\epsilon>0$.
    Como $f$ é contínua em $p$, existe $\delta_f>0$ tal que, para todo $x\in A$,
    \[
        \text{se }d(x,p)<\delta_f,\text{ então }d(f(x),f(p))<\epsilon.
    \]

    Como $\gamma$ é contínua em $t_0$ e $\gamma(t_0)=p$, existe $\delta_\gamma>0$ tal que, para todo $t\in I$,
    \[
        \text{se }|t-t_0|<\delta_\gamma,\text{ então }d(\gamma(t),p)<\delta_f.
    \]

    Assim, para todo $t\in I$ com $|t-t_0|<\delta_\gamma$,
    \[
        d(f(\gamma(t)),f(\gamma(t_0)))=d(f(\gamma(t)),f(p))<\epsilon.
    \]
    Portanto, $f\circ\gamma$ é contínua em $t_0$.

    A última afirmação é a contrapositiva da primeira.
\end{proof}

\begin{example}
    Seja $f:\mathbb R^2\setminus\{(0,0)\}\to\mathbb R$ dada por
    \[
        f(x,y)=\frac{xy}{x^2+y^2}.
    \]

    Vamos mostrar que $f$ não possui limite em $(0,0)$.

    Considere as curvas
    \[
        \gamma_1(t)=(t,0)
        \qquad\text{e}\qquad
        \gamma_2(t)=(t,t),
    \]
    definidas para $t\neq 0$.
    Ambas se aproximam de $(0,0)$ quando $t\to 0$.

    Ao longo de $\gamma_1$, temos
    \[
        f(\gamma_1(t))=f(t,0)=0.
    \]
    Logo,
    \[
        \lim_{t\to 0}f(\gamma_1(t))=0.
    \]

    Ao longo de $\gamma_2$, temos
    \[
        f(\gamma_2(t))=f(t,t)=\frac{t^2}{2t^2}=\frac{1}{2}.
    \]
    Logo,
    \[
        \lim_{t\to 0}f(\gamma_2(t))=\frac{1}{2}.
    \]

    Como obtivemos limites diferentes ao longo de duas curvas que se aproximam de $(0,0)$, concluímos que
    \[
        \lim_{(x,y)\to(0,0)}\frac{xy}{x^2+y^2}
    \]
    não existe.
\end{example}

\begin{example}
    Verificar apenas retas pode ser insuficiente.

    Defina $F:\mathbb R^2\to\mathbb R$ por
    \[
        F(x,y)=
        \begin{cases}
            \dfrac{x^2y}{x^4+y^2}, & \text{se }(x,y)\neq(0,0),\\
            0, & \text{se }(x,y)=(0,0).
        \end{cases}
    \]

    Ao longo de qualquer reta passando pela origem, obtemos limite $0$.
    De fato, se a reta é $y=ax$, então, para $a\neq 0$,
    \[
        F(t,at)
        =
        \frac{x^2y}{x^4+y^2}\bigg|_{(x,y)=(t,at)}
        =
        \frac{at^3}{t^4+a^2t^2}
        =
        \frac{at}{t^2+a^2}
        \longrightarrow 0.
    \]
    Se $a=0$, então $F(t,0)=0$.
    A reta vertical $x=0$ também dá $F(0,t)=0$.

    No entanto, ao longo da parábola $\gamma(t)=(t,t^2)$,
    \[
        F(\gamma(t))
        =
        F(t,t^2)
        =
        \frac{t^2t^2}{t^4+t^4}
        =
        \frac{1}{2}
    \]
    para todo $t\neq 0$.

    Portanto, $F$ não é contínua em $(0,0)$.
    Esse exemplo mostra que testar somente retas pode sugerir um comportamento falso.
\end{example}

O teste negativo por curvas é frequentemente o mais usado: encontramos uma curva ruim e concluímos que o limite não existe ou que a função não é contínua.
Mas existe também uma versão completa do teste, desde que aceitemos caminhos de teste cujos domínios sejam subconjuntos de $\mathbb R$, e não apenas intervalos.

\begin{definition}\index{caminho de teste}
    Seja $A\subseteq\mathbb R^n$ e seja $p\in\mathbb R^n$.

    Um \emph{caminho de teste em $A$ chegando a $p$} é uma função $\gamma:D\to A$, em que $D\subseteq\mathbb R$ possui $0$ como ponto de acumulação e
    \[
        \lim_{t\to 0}\gamma(t)=p.
    \]

    Dizemos que o caminho de teste é \emph{perfurado} se $\gamma(t)\neq p$ para todo $t\in D$.
\end{definition}

\begin{proposition}[Teste completo por caminhos]\index{limite!teste completo por caminhos}
    Seja $A\subseteq\mathbb R^n$, seja $f:A\to\mathbb R^m$, seja $p$ um ponto de acumulação de $A$ e seja $L\in\mathbb R^m$.

    Então $\lim_{x\to p}f(x)=L$ se, e somente se, para todo caminho de teste perfurado $\gamma:D\to A$ chegando a $p$, temos
    \[
        \lim_{t\to 0}f(\gamma(t))=L.
    \]
\end{proposition}

\begin{proof}
    A ida é a mesma ideia da proposição sobre limites ao longo de curvas.
    Se $\lim_{x\to p}f(x)=L$ e $\gamma:D\to A$ é um caminho de teste perfurado chegando a $p$, então, dado $\epsilon>0$, existe $\delta_f>0$ tal que
    \[
        x\in A\setminus\{p\},\ d(x,p)<\delta_f
        \quad\Longrightarrow\quad
        d(f(x),L)<\epsilon.
    \]
    Como $\gamma(t)\to p$, existe $\delta_\gamma>0$ tal que
    \[
        t\in D,\ 0<|t|<\delta_\gamma
        \quad\Longrightarrow\quad
        d(\gamma(t),p)<\delta_f.
    \]
    Como $\gamma$ é perfurado, $\gamma(t)\in A\setminus\{p\}$ para todo $t\in D$.
    Logo,
    \[
        d(f(\gamma(t)),L)<\epsilon
    \]
    sempre que $t\in D$ e $0<|t|<\delta_\gamma$.
    Portanto,
    \[
        \lim_{t\to 0}f(\gamma(t))=L.
    \]

    Provemos a recíproca por contraposição.
    Suponha que $L$ não é limite de $f$ em $p$.
    Então existe $\epsilon_0>0$ tal que, para todo $\delta>0$, podemos encontrar $x\in A\setminus\{p\}$ satisfazendo
    \[
        d(x,p)<\delta
        \qquad\text{e}\qquad
        d(f(x),L)\geq\epsilon_0.
    \]

    Para cada inteiro $k\geq 1$, aplique essa afirmação com $\delta=\frac{1}{k}$.
    Obtemos $x_k\in A\setminus\{p\}$ tal que
    \[
        d(x_k,p)<\frac{1}{k}
        \qquad\text{e}\qquad
        d(f(x_k),L)\geq\epsilon_0.
    \]

    Defina
    \[
        D=\left\{\frac{1}{k}:k\in\mathbb N,\ k\geq 1\right\}
    \]
    e $\gamma:D\to A$ por
    \[
        \gamma\left(\frac{1}{k}\right)=x_k.
    \]

    O conjunto $D$ possui $0$ como ponto de acumulação, e $\gamma(t)\to p$ quando $t\to 0$ em $D$, pois $d(x_k,p)<\frac{1}{k}$.
    Além disso, $\gamma$ é perfurado, pois $x_k\neq p$ para todo $k$.

    Porém,
    \[
        d(f(\gamma(1/k)),L)=d(f(x_k),L)\geq\epsilon_0
    \]
    para todo $k$.
    Portanto, $f\circ\gamma$ não tende a $L$ quando $t\to 0$ em $D$.

    Assim, a hipótese sobre todos os caminhos de teste perfurados implica necessariamente que
    \[
        \lim_{x\to p}f(x)=L.
    \]
\end{proof}

\begin{corollary}[Teste completo de continuidade]\index{continuidade!teste completo por caminhos}
    Seja $A\subseteq\mathbb R^n$, seja $f:A\to\mathbb R^m$ e seja $p\in A$.

    Então $f$ é contínua em $p$ se, e somente se, para todo conjunto $D\subseteq\mathbb R$ com $0\in D$ e $0$ ponto de acumulação de $D$, e para toda função $\gamma:D\to A$ contínua em $0$ tal que $\gamma(0)=p$, a função $f\circ\gamma:D\to\mathbb R^m$ é contínua em $0$.
\end{corollary}

\begin{proof}
    Se $p$ é ponto isolado de $A$, então $f$ é contínua em $p$.
    Além disso, se $\gamma:D\to A$ é contínua em $0$ e $\gamma(0)=p$, então, para $t$ suficientemente próximo de $0$, temos $\gamma(t)=p$.
    Logo, $f\circ\gamma$ é contínua em $0$.

    Suponha, então, que $p$ é ponto de acumulação de $A$.

    Se $f$ é contínua em $p$, então $f\circ\gamma$ é contínua em $0$ pela continuidade de composições.
    De fato, dado $\epsilon>0$, existe $\delta_f>0$ tal que
    \[
        x\in A,\ d(x,p)<\delta_f
        \quad\Longrightarrow\quad
        d(f(x),f(p))<\epsilon.
    \]
    Como $\gamma$ é contínua em $0$ e $\gamma(0)=p$, existe $\delta_\gamma>0$ tal que
    \[
        t\in D,\ |t|<\delta_\gamma
        \quad\Longrightarrow\quad
        d(\gamma(t),p)<\delta_f.
    \]
    Logo, para $t\in D$ com $|t|<\delta_\gamma$,
    \[
        d(f(\gamma(t)),f(\gamma(0)))=d(f(\gamma(t)),f(p))<\epsilon.
    \]
    Portanto, $f\circ\gamma$ é contínua em $0$.

    Reciprocamente, suponha que $f$ não é contínua em $p$.
    Como $p$ é ponto de acumulação de $A$, temos que $f(p)$ não é limite de $f$ em $p$.
    Pela prova da proposição anterior, existe um caminho de teste perfurado $\alpha:D_0\to A$ chegando a $p$ tal que $f(\alpha(t))$ não tende a $f(p)$ quando $t\to 0$ em $D_0$.

    Defina $D=D_0\cup\{0\}$ e $\gamma:D\to A$ por
    \[
        \gamma(0)=p
        \qquad\text{e}\qquad
        \gamma(t)=\alpha(t)\text{ se }t\in D_0.
    \]
    Então $\gamma$ é contínua em $0$ e $\gamma(0)=p$.
    Porém $f\circ\gamma$ não é contínua em $0$, pois não se aproxima de $f(p)=(f\circ\gamma)(0)$ ao longo dos pontos de $D_0$.

    Portanto, se $f\circ\gamma$ é contínua em $0$ para todo caminho de teste contínuo $\gamma$ com $\gamma(0)=p$, então $f$ é contínua em $p$.
\end{proof}

\section*{Exercícios}

\begin{exercise}
    Use curvas para mostrar que o limite abaixo não existe:
    \[
        \lim_{(x,y)\to(0,0)}\frac{x^2-y^2}{x^2+y^2}.
    \]
\end{exercise}

\begin{exercise}
    Seja $f:\mathbb R^2\to\mathbb R$ dada por
    \[
        f(x,y)=
        \begin{cases}
            \dfrac{y^2}{x^2+y^2}, & \text{se }(x,y)\neq(0,0),\\
            0, & \text{se }(x,y)=(0,0).
        \end{cases}
    \]
    Use curvas para decidir se $f$ é contínua em $(0,0)$.
\end{exercise}

\begin{exercise}
    Seja $g:\mathbb R^2\to\mathbb R$ dada por
    \[
        g(x,y)=
        \begin{cases}
            \dfrac{x^2y}{x^4+y^2}, & \text{se }(x,y)\neq(0,0),\\
            0, & \text{se }(x,y)=(0,0).
        \end{cases}
    \]
    Verifique que, ao longo de toda reta passando pela origem, o limite de $g$ é $0$.
    Em seguida, encontre uma curva que mostre que $g$ não é contínua em $(0,0)$.
\end{exercise}

\begin{exercise}
    Seja $F:\mathbb R^2\to\mathbb R^2$ dada por
    \[
        F(x,y)=
        \begin{cases}
            \left(\dfrac{xy}{x^2+y^2}, x+y\right), & \text{se }(x,y)\neq(0,0),\\
            (0,0), & \text{se }(x,y)=(0,0).
        \end{cases}
    \]
    Use curvas e o critério por coordenadas para mostrar que $F$ não é contínua em $(0,0)$.
\end{exercise}

\begin{exercise}
    Seja $f:A\to\mathbb R^m$ uma função contínua em $p\in A$, e seja $\gamma:I\to A$ uma curva contínua tal que $\gamma(t_0)=p$.
    Mostre diretamente pela definição de continuidade que $f\circ\gamma$ é contínua em $t_0$.
\end{exercise}

